%!TEX program = uplatex
\RequirePackage{plautopatch}
\documentclass[uplatex,dvipdfmx]{jlreq}
\usepackage{amsmath,amssymb}

\pagestyle{empty}

\begin{document}

%======================================================================
% 表紙
%======================================================================
\begin{center}
  \vspace*{20pt}
  {\Huge\bfseries トロイダル・コンパクト化}

  \vspace{30pt}

  {\Large\bfseries 対称領域の数論的商空間とその周辺}

  \vspace{50pt}

  {\large 前文}

  \vspace{10pt}

  \parbox{0.85\textwidth}{
    \large\bfseries
    数論的商空間のトロイダル・コンパクト化を中心に、
    トーリック多様体の入門から数論的構成まで幅広く論じる。
  }
\end{center}

\vfill

\newpage

%======================================================================
% 講演１：佐武一郎
%======================================================================
\noindent
{\large\bfseries トロイダル・コンパクト化について I, II}

\medskip
\hspace*{\fill}{\large\bfseries 佐武 一郎（中央大学・理工）}

\bigskip\bigskip

数論的商空間のコンパクト化として最初に考えられたもの（佐武・Baily--Borel のコンパクト化）は
標準的で最小という利点がある反面、非常に悪い特異点をもつものであった。
この欠点を除くために 1970 年代に井草、Hirzebruch 等により様々な非特異化が考えられたが、
その中で最も一般的かつ優れたものが、Mumford 達によって構成されたトロイダルコンパクト化である。

$\mathcal{D}$ を対称領域、$G$ をその自己同型群（の連結成分）、
$\Gamma$ をその（一つの $\mathbf{Q}$ 構造に関する）数論的部分群とするとき、
商空間 $\mathcal{V} = \Gamma \backslash \mathcal{D}$ を数論的商空間、
そのコンパクト化を数論的多様体という。
標準的コンパクト化 $\mathcal{V}^{\ast}$ は $\mathcal{V}$ に
$\mathcal{D}$ の有理境界成分 $\mathcal{F}$ から生じる低次元の数論的商空間
$\Gamma_{\mathcal{F}} \backslash \mathcal{F}$ を有限個つけ加えて得られる。
トロイダル・コンパクト化 $\bar{\mathcal{V}}^M$ は、さらに $\mathcal{V}^{\ast}$ を境界に沿って
blow up し、商空間 $\Gamma_{\mathcal{F}} \backslash \mathcal{F}$ を
その上のある種のファイバー空間で置き換えて得られるものである。
それは $\mathcal{D}$ の境界成分 $\mathcal{F}$ から生じる $\mathcal{V}$ の「穴」を埋めるために、
（$\mathcal{F}$ に関するジーゲル領域表示を使って）そのトーラス部分を
トーリック完備化で置き換えることによって構成される。

\medskip

このようにトロイダル・コンパクト化は対称領域の群論的構造をより精密に反映するものであり、
その境界上のファイバー空間の中にはアーベル多様体をファイバーとする空間（久賀のファイバー空間）、
さらにその上のトーリック束などが現れる。
したがってその構造を分析することは、これらファイバー空間およびそのコンパクト化
（いいかえれば、特異ファイバー）を調べる上に役立つと思われる。
またその上の line bundle を考察することによって、
種々の保型形式の間の関係（たとえば、Maass 形式と Jacobi 形式の関係）、
その幾何学的意味を解明できるかもしれない。
また Faltings による $\mathbf{Z}$ 上のスキームとしての構成が示すように、
数論的研究の対象としても興味あるものと思われる。

\medskip

例えていえば、この怪獣の代数幾何学的・数論的性質をもっと良く知る必要がある。
その健胆な胃の中には様々な動物達（上記のファイバー空間のコンパクト化、
そのファイバーであるアーベル多様体やトーリック束の退化したもの等）が飲み込まれている。
これら哀れな動物達を怪獣の腹中から助け出し、その本来の姿を
（より一般な枠組みの中で）回復させることは意味のある問題であろう。

\bigskip

第一の講演では、準備として、有界対称領域の境界成分、対応する放物的部分群の構造、
ジーゲル領域としての表示、典型的な例としてジーゲル空間の場合、等について概説する。

第二の講演では、数論的部分群 $\Gamma$ が与えられたとき、それに関する凸錐の分解、
トーリックな完備化、商空間 $\mathcal{V} = \Gamma \backslash \mathcal{D}$ の
$\mathcal{F}$ 方向の局所的完備化、それらの張り合わせによる全体的コンパクト化の構成等を
（主としてジーゲル空間の場合の例について）概説する。

\bigskip

この講演の補足として、三人の方々に講演をお願いすることにした。
第一日目に、トロイダル・コンパクト化の方法上また発想上、本質的なトーリック多様体について、
石井志保子さんに入門的な講演をお願いした。
第二日目にトロイダル・コンパクト化の重要な応用として、
上にも述べたアーベル多様体の退化の理論について、
その系統的研究を最初に手がけられた浪川幸彦氏に説明して頂くことにした。
最後に $\mathbf{Z}$ 上の構成については、この方面に詳しい藤原一宏氏に解説をお願いした。

\newpage

%======================================================================
% 講演２：石井志保子
%======================================================================
\noindent
{\large\bfseries トーリック多様体入門}

\medskip
\hspace*{\fill}{\large\bfseries 石井志保子（東工大・理）}

\bigskip\bigskip

$(\mathbf{C}^*)^n = \mathbf{C}^* \times \mathbf{C}^* \times \cdots \times \mathbf{C}^*$（$n$ 個の積）を $n$ 次元トーラスと呼ぶ。
これは積
\[
  (t_1, t_2, \ldots, t_n)(s_1, s_2, \ldots, s_n)
\]
により可換群になる。
いま代数多様体 $X$ にトーラス $(\mathbf{C}^*)^n$ が作用しているとする。
特に $X$ が正規で $(\mathbf{C}^*)^n$ と同型な開軌跡をもつとき、$X$ を $n$ 次元トーリック多様体と呼ぶ。

このトーリック多様体はとても把握しやすい多様体なので、
代数多様体の良いモデルとして、また新しい多様体を具体的に構成する場合にしばしば用いられる。
ではなぜ把握しやすいのか、どうやって把握するのかを紹介するのがこの講演の目的である。

講演では $n$ 次元実線形空間の上に扇と呼ばれる有理多面錐の族を考える。
この $n$ 次元実線形空間の上の扇と $n$ 次元トーリック多様体とは一対一に対応することを紹介する。
そしてトーリック多様体の色々な性質は扇の性質で記述されることを見る。
したがってトーリック多様体は扇によって完全に把握されるということになる。

\newpage

%======================================================================
% 講演３：藤原一宏
%======================================================================
\noindent
{\large\bfseries 数論的トロイダルコンパクト化}

\medskip
\hspace*{\fill}{\large\bfseries 藤原 一宏（名大・多元数理）}

\bigskip\bigskip

有界対称領域の算術商をみると、数学者は胸がときめく。
この講演では数論の人間のときめきを語ってみたい。
まず、予備知識をいくつか。
そのような多様体にはある標準的な仕方で代数体上の代数多様体の構造が入ることが、
志村五郎氏（プリンストン大）の一連の仕事、及びそれを補完する多くの人たちの仕事により
示されている（それ故に志村多様体といわれる）。
この事実は現在の数論において極めて重要な役割をはたしている。

まず、これが虚数乗法論の一般化（の一つの方向）であること、
つまりモデュラー関数の特殊値がどのようなアーベル拡大を生成するかに解答を与えている
（Kronecker の青春の夢、Hilbert の第12問題）。
また、そのエタールコホモロジー群上にはガロア群とアデール化された代数群が共に作用し、
ガロア表現と保型表現との間の非可換類体論を実現すると信じられている
（コホモロジーの研究のためにもコンパクト化が重要なのである）。

\medskip

志村多様体自身の話題に戻ろう。
志村氏の仕事にもアーベル多様体という形で現れているが、
Deligne は Grothendieck のモチーフの哲学
（モチーフという数学的実在を信じる宗教のようなもの）と志村氏の仕事を結びつけた。
夢を語るとすれば、志村多様体は（条件を付けた）モチーフのモデュライ空間なのである。
（有界対称領域自身は Hodge 構造の分類空間と解釈される。「モチーフは Hodge 構造を生む」のである。）

このようなモチーフ論的視点に立つと、
トロイダルコンパクト化の構成はモチーフの退化理論の帰結としてとらえられることになる
（残念ながら、整数環 $\mathbf{Z}$ 上満足のいく形でのモチーフ理論は知られていないのだが）。
ここでは特殊なモチーフであるアーベル多様体とその退化である $1$-motive を通して
PEL 型の志村多様体のトロイダルコンパクト化が $\mathbf{Z}$ 上自然に構成できることを述べたい
（Siegel modular 多様体の場合は Chai 及び Faltings により、PEL 型の時は筆者による）。
しかしながら、あまり一般的には行わず、
楕円モデュラー、ヒルベルトモデュラーなどのより親しみやすい例から入りたいと思う。

\end{document}
